\insertImage{c4a662b0-2b67-4605-b38f-48b9e903fd03.png}{Aquarelle représentant deux salles de classe symbolisant la transition entre les systèmes éducatifs traditionnels et modernes.}

\chapter{Système Éducatif et Formation}

Faites ce petit test : repensez à votre scolarité, et cherchez le moment où on vous a appris à travailler sans qu'on vous dise quoi faire.

Prenez votre temps. Pour la plupart d'entre nous, la réponse est : jamais. De la maternelle au diplôme, nous avons appris à répondre à des consignes, à être notés par une autorité, à lever la main pour parler. Ce sont peut-être des compétences utiles. Mais avouez que, comme préparation à l'entreprise libérée, on a connu mieux.

Ce chapitre explore ce paradoxe fondateur : nous demandons à des adultes de s'autogérer après leur avoir enseigné pendant vingt ans à obéir. Et devinez qui doit combler l'écart ? L'entreprise. C'est un travail considérable, presque toujours sous-estimé — et son omission est l'une des causes d'échec les plus silencieuses du modèle.

\section{La formation traditionnelle en management et leadership}

Le problème commence tôt et se prolonge haut : jusque dans la formation de ceux qui dirigent.

Les sociologues de l'éducation ont montré depuis longtemps que l'école française ne transmet pas que des savoirs : elle transmet un rapport à l'autorité et une hiérarchie des positions, qu'elle contribue à reproduire de génération en génération\footnote{Bourdieu, P., \& Passeron, J.-C. (1970). La Reproduction. Éléments pour une théorie du système d'enseignement. Les Éditions de Minuit.}. Le système des grandes écoles pousse cette logique à son sommet : sélection précoce, classement permanent, et la promesse implicite que l'excellence académique fonde la légitimité à commander — l'exact contraire de la légitimité par la confiance que suppose une entreprise libérée.

Quant à la formation au management proprement dite, elle a ses propres critiques, et pas des tendres. Henry Mintzberg a consacré un livre entier à démontrer que les écoles de gestion enseignent l'analyse et le contrôle à des gens qui n'ont jamais managé personne — et forment ainsi des dirigeants habiles en tableurs et démunis face à l'humain\footnote{Mintzberg, H. (2004). Managers Not MBAs: A Hard Look at the Soft Practice of Managing and Management Development. Berrett-Koehler Publishers.}. En France, le sociologue François Dupuy dresse un constat voisin : une pensée managériale enfermée dans les process et le reporting, qui confond contrôler et diriger — et qui se reproduit d'autant mieux qu'elle rassure\footnote{Dupuy, F. (2015). La faillite de la pensée managériale: Lost in management 2. Seuil.}.

Résultat : le manager français type arrive aux responsabilités avec un logiciel mental — planifier, organiser, contrôler — que la libération lui demande précisément de désinstaller. On comprend mieux pourquoi tant de transformations butent sur l'encadrement intermédiaire : on ne lui a jamais appris autre chose.

\section{L'impact de l'éducation sur l'adoption de nouveaux modèles de gestion}

Ce conditionnement éducatif produit des effets très concrets au moment d'adopter un modèle libéré — et pas seulement chez les managers.

Chez les dirigeants, il façonne ce qui est pensable : quand toute votre formation a présenté l'organisation comme une machine à optimiser, une entreprise fondée sur la confiance ressemble à de la littérature. Les institutions internationales elles-mêmes soulignent l'écart croissant entre ce que le monde du travail exige — autonomie, coopération, apprentissage permanent — et ce que les systèmes de formation produisent\footnote{OCDE (2019). OECD Employment Outlook 2019: The Future of Work. OECD Publishing.}.

Chez les salariés, le conditionnement joue plus subtilement. Vingt ans d'école laissent des réflexes : attendre la consigne, craindre l'erreur, viser la note. Mettez ces réflexes dans une entreprise qui dit « prends l'initiative, l'erreur est permise, il n'y a plus de note », et vous obtenez la scène que tous les libérateurs racontent : la liberté offerte, et personne ne s'en saisit pendant des mois. Ce n'est ni de la paresse ni du cynisme. C'est de la prudence apprise — profondément apprise.

D'où une conséquence stratégique que les entreprises qui réussissent ont comprise : la formation n'est pas un accompagnement de la libération, c'en est une condition préalable. Michelin, dans sa démarche de responsabilisation, a investi massivement dans la préparation de ses équipes et de ses managers avant de leur transférer les décisions\footnote{Getz, I. (2017). L'entreprise libérée. Fayard.}. Transférer le pouvoir sans transférer les compétences, ce n'est pas de la confiance — c'est de l'abandon.

\section{L'éducation et la formation nécessaires pour soutenir une entreprise libérée}

Alors, que faudrait-il apprendre — et désapprendre ? Le référentiel tient en trois blocs, et vous remarquerez qu'aucun ne figure au programme de nos écoles.

Pour ceux qui dirigent, le cœur est ce que Robert Greenleaf appelait le leadership serviteur : concevoir son rôle comme étant au service de ceux qu'on dirige, et non l'inverse\footnote{Greenleaf, R. K. (1977). Servant leadership: A journey into the nature of legitimate power and greatness. Paulist Press.}. Ça s'apprend — mais pas en cours magistral : par la pratique accompagnée, le feedback, et une bonne dose de travail sur soi. Pour un dirigeant formé à la verticalité, c'est moins une formation qu'une rééducation, au sens presque kinésithérapique du terme : des réflexes à défaire un par un.

Pour tout le monde, deux familles de compétences deviennent critiques. Les compétences relationnelles, d'abord — écouter vraiment, formuler un désaccord, recevoir une critique — qui cessent d'être du « savoir-être » décoratif pour devenir l'infrastructure même de la coordination\footnote{Goulston, M. (2009). Just Listen: Discover the Secret to Getting Through to Absolutely Anyone. Amacom.}. Les compétences de pilotage de soi, ensuite : fixer ses priorités, donner et demander du feedback en continu plutôt qu'attendre une évaluation annuelle — une évolution que même les entreprises classiques ont entamée\footnote{Cappelli, P., \& Tavis, A. (2016). The Performance Management Revolution. Harvard Business Review, 94(10), 58-67.}.

Le dernier bloc est le plus profond, et Warren Bennis l'avait formulé bien avant la mode des entreprises libérées : devenir leader, c'est devenir soi-même — un travail d'expression personnelle plus que d'acquisition technique\footnote{Bennis, W. (2009). On Becoming a Leader. Basic Books.}. Une entreprise libérée est exigeante en maturité humaine. Cette maturité se cultive ; encore faut-il que quelqu'un s'en charge, puisque ni l'école ni l'université ne s'y sont attelées.

De ma propre expérience de dirigeant en holacratie, je retiens un programme minimal, que j'appliquerais dès le premier jour si c'était à refaire. Former chacun à deux choses avant de lui transférer les décisions. D'abord, la vie économique d'une entreprise : ce que coûte réellement un salaire, ce qu'exige un client, ce qu'imposent la loi et la trésorerie — parce qu'on ne peut pas décider librement dans un monde dont on ignore les règles. Ensuite, le fonctionnement du modèle lui-même — ses rituels, ses rôles, ses limites. À partir de là, toute l'autonomie nécessaire. Avant, aucune : ce ne serait pas de la liberté, ce serait un piège.

Car la formation est le chaînon manquant entre l'intention et la réussite. Quand une libération échoue, on accuse le modèle, le dirigeant, la culture. On devrait plus souvent se demander qui, dans cette histoire, avait appris à faire ce qu'on lui demandait soudain de faire.
